% !TeX program = pdflatex
% Samodzielny skrypt: Teoria kategorii — podstawy i przyklady.
\documentclass[11pt,a4paper]{article}
\usepackage[T1]{fontenc}
\usepackage[utf8]{inputenc}
\usepackage[polish]{babel}
\usepackage{lmodern}
\usepackage{amsmath,amssymb,amsthm,mathtools,amscd}
\usepackage{mathrsfs}
\usepackage{microtype}
\usepackage{xcolor}
\usepackage{geometry}
\usepackage{enumitem}
\usepackage[unicode]{hyperref}
\hypersetup{colorlinks=true,linkcolor=blue!45!black,urlcolor=blue!45!black}
\geometry{margin=2.5cm}
\setlength{\parindent}{0pt}
\setlength{\parskip}{4pt}
\setlength{\emergencystretch}{2em}
\setlist[enumerate]{topsep=4pt,itemsep=2pt,leftmargin=*}
\setlist[itemize]{topsep=4pt,itemsep=2pt,leftmargin=*}
\numberwithin{equation}{section}
\theoremstyle{plain}
\newtheorem{twierdzenie}{Twierdzenie}[section]
\newtheorem{stwierdzenie}[twierdzenie]{Stwierdzenie}
\newtheorem{lemat}[twierdzenie]{Lemat}
\newtheorem{wniosek}[twierdzenie]{Wniosek}
\theoremstyle{definition}
\newtheorem{definicja}[twierdzenie]{Definicja}
\newtheorem{przyklad}[twierdzenie]{Przykład}
\newtheorem{cwiczenie}{Ćwiczenie}
\theoremstyle{remark}
\newtheorem{uwaga}[twierdzenie]{Uwaga}
\newcommand{\cat}[1]{\mathsf{#1}}
\newcommand{\Set}{\cat{Set}}
\newcommand{\Grp}{\cat{Grp}}
\newcommand{\Mon}{\cat{Mon}}
\newcommand{\Ring}{\cat{Ring}}
\newcommand{\Top}{\cat{Top}}
\newcommand{\Vect}{\cat{Vect}}
\newcommand{\Ab}{\cat{Ab}}
\newcommand{\Pos}{\cat{Pos}}
\newcommand{\Rel}{\cat{Rel}}
\newcommand{\id}{\mathrm{id}}
\newcommand{\Hom}{\operatorname{Hom}}
\newcommand{\Ob}{\operatorname{Ob}}
\newcommand{\Nat}{\operatorname{Nat}}
\newcommand{\op}{\mathrm{op}}
\newcommand{\im}{\operatorname{im}}
\newcommand{\Ker}{\operatorname{ker}}
\newcommand{\colim}{\operatorname{colim}}
\DeclareMathOperator{\End}{End}
\DeclareMathOperator{\Aut}{Aut}
\DeclareMathOperator{\ev}{ev}
\title{\textbf{Teoria kategorii od podstaw}\\\large Zbiór ćwiczeń do samodzielnego rozwiązania z podpowiedziami}
\author{}
\date{}
\begin{document}
\maketitle
\section{Ćwiczenia do samodzielnego rozwiązania}\label{sec:cwiczenia}

Zbiór towarzyszy skryptowi \emph{Teoria kategorii od podstaw}. Zawiera wyłącznie zadania przeznaczone do samodzielnego rozwiązania oraz krótkie wskazówki, bez gotowych rozwiązań. Zachowano zakres i kolejność ćwiczeń z pierwotnej wersji skryptu.

\begin{cwiczenie}[Kategoria z dwoma obiektami]
Zbuduj kategorię o obiektach $0,1$, w której istnieje dokładnie jeden nietrywialny morfizm $0\to1$ i nie istnieje morfizm $1\to0$. Wypisz wszystkie możliwe złożenia i wskaż obiekt początkowy oraz końcowy.

\emph{Wskazówka:} jest to kategoria porządku $0\leq1$.
\end{cwiczenie}

\begin{cwiczenie}[Izomorfizmy w kategorii monoidu]
Wykaż, że w kategorii odpowiadającej monoidowi $M$ izomorfizmy są dokładnie elementami odwracalnymi monoidu. Co otrzymujemy dla monoidu $(\mathbb N,+,0)$?

\emph{Wskazówka:} równości $g\circ f=\id$ i $f\circ g=\id$ stają się równościami w monoidzie.
\end{cwiczenie}

\begin{cwiczenie}[Przestrzenie liniowe]
Udowodnij, że w $\Vect_k$ monomorfizmy to dokładnie odwzorowania liniowe różnowartościowe, a epimorfizmy to dokładnie odwzorowania liniowe surjektywne.

\emph{Wskazówka:} w części dotyczącej monomorfizmów wykorzystaj $k\to V$, $1\mapsto v$; w części dotyczącej epimorfizmów rozważ rzut $W\to W/\im f$.
\end{cwiczenie}

\begin{cwiczenie}[Mono i epi bez izomorfizmu]
W kategorii posetu $\{0<1\}$ sprawdź, że jedyna nietrywialna strzałka jest jednocześnie monomorfizmem i epimorfizmem, ale nie jest izomorfizmem.

\emph{Wskazówka:} hom-zbiory mają co najwyżej jeden element.
\end{cwiczenie}

\begin{cwiczenie}[Iloczyn w porządku]
Rozważ poset dzielników liczby $30$ uporządkowany przez podzielność. Oblicz produkty i koprodukty par $(2,3)$, $(6,10)$ i $(5,15)$.

\emph{Wskazówka:} iloczyn odpowiada $\gcd$, a koprodukt $\operatorname{lcm}$.
\end{cwiczenie}

\begin{cwiczenie}[Iloczyn i koprodukt przestrzeni liniowych]
Niech $V,W$ będą skończenie wymiarowymi przestrzeniami nad $k$. Wykaż z własności uniwersalnej, że $V\oplus W$ jest produktem w $\Vect_k$. Dlaczego $V\oplus W$ jest także koproduktem?

\emph{Wskazówka:} porównaj wzory $x\mapsto(f(x),g(x))$ oraz $(v,w)\mapsto a(v)+b(w)$.
\end{cwiczenie}

\begin{cwiczenie}[Kontrawariancja]
Dla funkcji $f:X\to Y$ sprawdź równości $f^{-1}(B_1\cup B_2)=f^{-1}(B_1)\cup f^{-1}(B_2)$ i $f^{-1}(Y\setminus B)=X\setminus f^{-1}(B)$. Dlaczego przyporządkowanie $X\mapsto\mathcal P(X)$, $f\mapsto f^{-1}$ definiuje funktor kontrawariantny?

\emph{Wskazówka:} porównaj kierunki odwzorowań oraz równość $(g\circ f)^{-1}=f^{-1}\circ g^{-1}$.
\end{cwiczenie}

\begin{cwiczenie}[Naturalność nie zależy od współrzędnych]
Rozpisz naturalność $\iota_V:V\to V^{**}$ dla odwzorowania liniowego $f:V\to W$. Wskaż dokładnie, w którym miejscu używasz definicji $f^*$.

\emph{Wskazówka:} oblicz obie strony na $v\in V$ i $\mu\in W^*$.
\end{cwiczenie}

\begin{cwiczenie}[Element uniwersalny]
Niech $\mathcal C$ będzie kategorią lokalnie małą, $F:\mathcal C\to\Set$ funktorem, $A$ obiektem kategorii $\mathcal C$ oraz $u\in F(A)$. Załóż, że dla każdego obiektu $X$ przyporządkowanie $f\mapsto F(f)(u)$ jest bijekcją $\Hom_{\mathcal C}(A,X)\to F(X)$. Udowodnij, że jeżeli para $(A',u')$ ma tę samą własność, to istnieje dokładnie jeden izomorfizm $\theta:A\to A'$ taki, że $F(\theta)(u)=u'$.

\emph{Wskazówka:} zastosuj uniwersalność najpierw do $u'\in F(A')$, a następnie w odwrotnym kierunku.
\end{cwiczenie}

\begin{cwiczenie}[Włókna]
Niech $f:\mathbb Z\to\mathbb Z/3\mathbb Z$ będzie redukcją modulo $3$. Opisz pullback $\mathbb Z\times_{\mathbb Z/3\mathbb Z}\{*\}$, gdzie $*\mapsto[1]$.

\emph{Wskazówka:} określ pary $(n,*)$, dla których $[n]=[1]$.
\end{cwiczenie}

\begin{cwiczenie}[Sklejanie]
Skonstruuj pushout w $\Set$ dwóch funkcji $f:C\to A$, $g:C\to B$, gdy $C=\{c_1,c_2\}$, $A=\{a_1,a_2\}$, $B=\{b_1,b_2\}$, $f(c_i)=a_i$, $g(c_1)=g(c_2)=b_1$.

\emph{Wskazówka:} opisz klasy relacji równoważności na $A\sqcup B$.
\end{cwiczenie}

\begin{cwiczenie}[Yoneda w $\Set$]
Wyznacz wszystkie transformacje naturalne z funktora $X\mapsto\Hom_{\Set}(\{0,1\},X)$ do funktora $\id_{\Set}$. Opisz je bez korzystania z lematu Yonedy, a następnie sprawdź zgodność wyniku z jego tezą.

\emph{Wskazówka:} powinny istnieć dokładnie dwie: wybór pierwszej lub drugiej współrzędnej.
\end{cwiczenie}

\begin{cwiczenie}[W pełni wierny funktor]
Sprawdź bezpośrednio, że funktor $\cat{Mat}_k\to\Vect_k^{\mathrm{fd}}$ jest w pełni wierny. Dlaczego nie jest on izomorfizmem kategorii?

\emph{Wskazówka:} pierwsza część to reprezentacja odwzorowania liniowego macierzą; druga wymaga porównania obiektów obu kategorii.
\end{cwiczenie}

\begin{cwiczenie}[Wolna przestrzeń]
Dla zbioru $S=\{a,b\}$, ciała $k=\mathbb R$ i funkcji $\varphi:S\to\mathbb R^2$ zadanej przez $\varphi(a)=(1,2)$, $\varphi(b)=(-1,3)$ oblicz macierz jedynego odwzorowania liniowego $\widetilde\varphi:\mathbb R^{(S)}\to\mathbb R^2$ w bazach $(e_a,e_b)$ i standardowej.

\emph{Wskazówka:} kolumnami macierzy są obrazy generatorów.
\end{cwiczenie}

\begin{cwiczenie}[Tożsamości trójkątne]
Sprawdź obie tożsamości trójkątne dla sprzężenia funktora przestrzeni dyskretnej $D:\Set\to\Top$ z funktorem zapominania $U:\Top\to\Set$.

\emph{Wskazówka:} jednostka jest identycznością na zbiorach, a kojednostka $D(U(X))\to X$ jest identycznością jako funkcja, lecz niekoniecznie homeomorfizmem.
\end{cwiczenie}

\begin{cwiczenie}[Granica odwrotna]
Dla każdego $n\geq0$ niech $X_n=\{0,1,\ldots,n\}$ oraz $p_n:X_{n+1}\to X_n$, $p_n(x)=\min(x,n)$. Opisz wszystkie zgodne ciągi $(x_n)_{n\geq0}$ i ich granicę odwrotną.

\emph{Wskazówka:} rozważ ciągi odpowiadające ustalonej liczbie naturalnej oraz jeden ciąg ``rosnący bez ograniczenia''.
\end{cwiczenie}

\begin{cwiczenie}[Lista jako monada]
Dla $X=\{a,b\}$ oblicz $\mu_X([[a,b],[],[a]])$ oraz oba złożenia występujące w aksjomacie łączności dla listy $[[[a],[b]],[[a,b]]]$ elementów $T^3(X)$.

\emph{Wskazówka:} w obu przypadkach po całkowitym spłaszczeniu otrzymasz tę samą listę.
\end{cwiczenie}


\begin{cwiczenie}[Konstrukcja granicy w zbiorach]
Niech $D$ będzie diagramem w $\Set$ o obiektach $A=\{1,2,3\}$, $B=\{0,1\}$, $C=\{u,v\}$ i strzałkach $f:A\to C$, $g:B\to C$, gdzie $f(1)=f(2)=u$, $f(3)=v$, $g(0)=u$, $g(1)=v$. Wypisz wszystkie elementy pullbacku $A\times_C B$ i zweryfikuj własność uniwersalną dla $X=\{*\}$.

\emph{Wskazówka:} znajdź pary $(a,b)$ spełniające $f(a)=g(b)$; sprawdzaj morfizmy z $X$ do obiektu produktu włóknistego osobno dla każdego możliwego elementu.
\end{cwiczenie}

\begin{cwiczenie}[Naturalne operacje na krotkach]
Dla liczb naturalnych $m,n\geq1$ wyznacz liczbę transformacji naturalnych między funktorami $X\mapsto X^n$ i $X\mapsto X^m$ na $\Set$. Zastosuj wynik dla $n=3$, $m=2$ i wypisz wszystkie operacje.

\emph{Wskazówka:} zastosuj lemat Yonedy do obiektu $\{1,\ldots,n\}$, a następnie licz funkcje $\{1,\ldots,m\}\to\{1,\ldots,n\}$.
\end{cwiczenie}

\begin{cwiczenie}[Sprzężenie kwantyfikatorów]
Dla funkcji $f:\{1,2,3,4\}\to\{a,b\}$ określonej przez $f(1)=f(2)=a$ i $f(3)=f(4)=b$ oblicz $\exists_f(A)$ oraz $\forall_f(A)$ dla $A=\{1,2,3\}$. Zdefiniuj $\exists_f(A)=f(A)$ oraz $\forall_f(A)=\{y\in\{a,b\}:f^{-1}(\{y\})\subseteq A\}$. Sprawdź dla $B=\{a\}$ obie równoważności
$\exists_f(A)\subseteq B\iff A\subseteq f^{-1}(B)$ i
$f^{-1}(B)\subseteq A\iff B\subseteq\forall_f(A)$.

\emph{Wskazówka:} zbiór $\forall_f(A)$ zawiera tylko te wartości, których \emph{całe} włókno należy do $A$.
\end{cwiczenie}

\begin{cwiczenie}[Koequalizator jako iloraz grupy]
Dla homomorfizmów $f,g:\mathbb Z\to\mathbb Z$ danych przez $f(n)=2n$, $g(n)=5n$ opisz ich koequalizator w $\Ab$ i udowodnij bezpośrednio jego własność uniwersalną.

\emph{Wskazówka:} rozważ iloraz $\mathbb Z/\im(f-g)$; określ najpierw obraz $f-g$.
\end{cwiczenie}

\begin{cwiczenie}[Dwie różne monady]
Dla $X=\{a,b\}$ porównaj mnożenia monady podzbiorów i monady list, stosując je odpowiednio do $\{\{a,b\},\{a\}\}$ oraz $[[a,b],[a]]$. Wyjaśnij, jakie informacje znikały po każdym mnożeniu.

\emph{Wskazówka:} lista pamięta kolejność i wielokrotność wystąpień elementów; zbiór ich nie pamięta.
\end{cwiczenie}


\end{document}
